\section*{Chapter \ref{ch:ml:monitoring-and-retraining} --- Monitoring and Retraining}

\begin{solution}{exo:ml:monitoring-and-retraining:1}
Inputs: the slow fade's growing spread (seen by the feature index and Kolmogorov--Smirnov statistic, by nothing else). Outputs:
the unit change as a rise of the share inside the trading threshold from 50\% to 63\%, which the inputs also see but which only
the outputs translate into trading (a quarter fewer names traded). Outcomes: the reversal, invisible to inputs and outputs because
neither the features nor the distribution of predictions change.
\end{solution}

\begin{solution}{exo:ml:monitoring-and-retraining:2}
With 200 names, each decile holds 20 names, and sampling noise alone gives each feature's daily index a spread that puts the
largest of five above 0.10 on more than one day in twenty. The rule of thumb was set for large credit portfolios and says nothing
about error rates; a threshold of 0.10 here would page more often than once a month, and the team would learn to ignore it.
\end{solution}

\begin{solution}{exo:ml:monitoring-and-retraining:3}
Its false pages arrive at random, about one day in 21, so the wait for the first one after any date is geometric with a median of
about 14 days (the table's ``none'' column shows 10 to 17 for the five monitors). A first alarm means something only when it
comes well before that.
\end{solution}

\begin{solution}{exo:ml:monitoring-and-retraining:4}
$\Delta=0.103-0.087=0.017$ (unrounded 0.0168), $\sigma=0.071$, $z_{0.95}+z_{0.90}=1.645+1.282=2.927$: $n=(2.927\times0.071/0.0168)^2
\approx150$ days. The daily coefficient's standard deviation shrinks roughly as one over the square root of the number of names;
with 2,000 names it is about $0.071/\sqrt{10}=0.022$, and $n$ falls tenfold to about 15 days.
\end{solution}

\begin{solution}{exo:ml:monitoring-and-retraining:5}
The test's standard error is estimated from the days seen so far. On three days the sample standard deviation of the differences
can be tiny by chance, the estimated standard error with it, and a modest mean difference then looks overwhelming: with verdicts
from the third day, one of the chapter's runs shows a p-value of $1.3\times10^{-6}$ on that day. The test's error guarantee
assumes a known standard error; ten days are a burn-in that makes the estimate usable. Without a failure, verdicts from the third
day promote 16 of 100 null challengers, against 2 with the burn-in.
\end{solution}

\begin{solution}{exo:ml:monitoring-and-retraining:6}
The alarm was right and the fix was wrong. The vendor's change was a data error; retraining taught the model to multiply the
feature in lots by a hundred times the coefficient, burying the error in the model. When the vendor corrects the feed, the
retrained model is wrong by a factor of a hundred on that feature, and the alarm fires again. Fix the data (convert units at the
boundary, add a schema or range check), keep the model.
\end{solution}

\begin{solution}{exo:ml:monitoring-and-retraining:7}
\texttt{ml\_monitor.pooled}: the pooled index's threshold at one page a month is 0.0056, twenty times below the daily one, since
4,000 samples a feature carry much less noise than 200. Consecutive windows share 19 of their 20 days, so the statistic moves
slowly and its false alarms come in spells, 3.7 alarm days per page on average. The first alarm does not beat a blind monitor
(median 15.5 days for the fade, 18 without a failure), and the five-in-21 rule is tripped by a single false spell (median 20
days without a failure), so neither count means what it did. Measured by a month of continuous alarm, the pooled index sees the
fade after a median of 22.5 days, against 120 for the daily index; but four of the 20 runs without failures also show a false
month of continuous alarm. Pooling buys sensitivity to slow drifts with alarms that are fewer but longer, and harder to dismiss.
\end{solution}

\begin{solution}{exo:ml:monitoring-and-retraining:8}
No input or output monitor can see it: the features' distribution is unchanged, and so, to first order, is that of the
predictions. Only outcomes can, and the loss is gradual: the coefficient falls to 0.081 at the end of the fade, and the chapter's
formula needs about 80 days at that level, so the alarm would come months after the fade began, if at all at one false page a
month. The practical answer is not a monitor but a schedule: refit or revalidate the model regularly (chapter~12), compare each
feature's recent contribution (its coefficient refitted on a rolling window) with the training one, and look at attribution of
the realised IC by feature, which concentrates the evidence on the feature that fades.
\end{solution}

\begin{solution}{pb:ml:monitoring-and-retraining:1}
\textbf{Part I.}
\begin{enumerate}
\item Two hundred names a day, five standardised features, next-day returns loading 0.08, $-0.06$ and 0.05 on the first three
plus unit noise; a linear model fitted on 60 days runs for 400, trading the half of names with the largest forecasts. From day
200: the second feature a hundred times smaller; the first feature's effect reversed; the third's effect fading to zero over 150
days while its spread grows by half.
\item Inputs (feature distributions against the training reference; at once), outputs (distribution of predictions, share inside
the trading threshold; at once), outcomes (realised information coefficient, P\&L; after the horizon, and noisy).
\item $\sum_b(q_b-r_b)\ln(q_b/r_b)$ over the reference's deciles; thresholds 0.115 (largest feature index), 0.115 (largest
Kolmogorov--Smirnov statistic), 0.089 (index of the predictions), 0.07 (change in the share inside the threshold).
\item Because a person is paged once per episode, not once per day of it; calibrating on alarm days would make persistent
statistics look noisier than they are.
\end{enumerate}
\textbf{Part II.}
\begin{enumerate}[resume]
\item One-sided, on the daily coefficient's shortfall below its expected 0.103, allowance $k=0.035$ (half the daily standard
deviation of 0.071), threshold $h=0.11$ chosen on runs without failures to page once a month, reset after each page.
\item About $((z_{0.95}+z_{0.90})\sigma/\Delta)^2$: 4 days for the reversal, about 150 for the unit change, about 80 for the fade at its
end.
\item The unit change removes one feature's contribution from the forecast; the coefficient falls only from 0.103 to 0.087, a
drop a quarter of a day's noise.
\item A first alarm is the first page after the failure, which a blind monitor also produces within about two weeks; a sustained
alarm is five alarm days in 21, which no monitor reaches without a failure.
\end{enumerate}
\textbf{Part III.}
\begin{enumerate}[resume]
\item All four input and output monitors on the first day, sustained at once; the outcome monitor pages once after a median of
9.5 days and never insists.
\item Only the IC CUSUM: within two days in 13 of 20 runs, sustained at once.
\item Nothing sees its lost accuracy; the input monitors see its growing spread, sustained after a median of 42 days (feature
index) and 60 (Kolmogorov--Smirnov).
\item Nothing in the table; a scheduled revalidation, or a monitor on each feature's contribution to the realised coefficient.
\end{enumerate}
\textbf{Part IV.}
\begin{enumerate}[resume]
\item \emph{The Tuesday the model went quiet.} At one false page a month (medians over 20 runs; first alarm / sustained):
unit change, every input and output monitor 0 / 0 days, IC CUSUM 9.5 days / never; reversal, input and output monitors no better
than blind (10 to 17 days, never sustained), IC CUSUM 1 / 0 days; slow fade, feature index 12 / 42 days, Kolmogorov--Smirnov 8 /
60, prediction index, share inside and IC CUSUM never sustained.
\item A challenger refitted on 40 post-reversal days (information coefficient 0.105 against the champion's 0.003) is promoted in
all 20 runs after a median of 10 shadow days, at most 19; without a failure, 2 of 100 challengers are promoted.
\item A challenger with a sign error is rolled back by the CUSUM after a median of 3 days; the information coefficient given up
is 0.26 IC-days at full size and 0.026 at a tenth.
\item Page the data owner on the share-inside and feature alarms, stop or shrink the model, find the unit change and fix the feed;
not retrain.
\item Retrain after a confirmed change in the relation, once there is enough post-change data, and on a schedule; retire when
refitted challengers no longer beat a null.
\item Input alarms: the data owner; output alarms: the desk and model owner; outcome alarms: the model owner and risk.
\item Every alarm, who acknowledged it, the action (data fix, size cut, rollback, promotion) with its reason, linked to the model
version in the registry.
\item A false-alarm rate the team will answer and a measured delay on the failures it is for.
\end{enumerate}
\end{solution}

\begin{solution}{iq:ml:monitoring-and-retraining:1}
Inputs (schema, freshness, feature distributions), outputs (prediction distribution, positions, turnover, share of trades) and
outcomes (realised IC, P\&L, mark-outs). Each layer sees failures the others miss: a unit change is loud in the inputs and faint in
the outcomes; a regime change is silent in the inputs and loud in the outcomes.
\iqlookfor{three layers and an example each.}
\end{solution}

\begin{solution}{iq:ml:monitoring-and-retraining:2}
$\sum_b(q_b-r_b)\ln(q_b/r_b)$ over the reference's bins. Choose the threshold for a false-alarm rate on data without failures
(history or simulation), with the actual sample size and number of features monitored; the 0.10 and 0.25 rules of thumb ignore
both.
\iqlookfor{the formula and calibration to an error rate.}
\end{solution}

\begin{solution}{iq:ml:monitoring-and-retraining:3}
$\Delta=0.02$, $\sigma=0.07$: $((1.645+1.282)\times0.07/0.02)^2\approx105$ days, five months, for a one-sided 5\% test with 90\%
power; more names shorten it, as $\sigma$ falls roughly with the square root of their number.
\iqlookfor{the sample-size arithmetic and its dependence on breadth.}
\end{solution}

\begin{solution}{iq:ml:monitoring-and-retraining:4}
Shadow: the new model sees live inputs and its decisions are scored but not executed; right for comparing forecasts at no risk.
Canary: it trades a small share of capital; right when its trading changes outcomes (impact, fills). \omterm{def:ml:monitoring-and-retraining:shadow}{Champion--challenger}: the
discipline around both, keeping the incumbent until the challenger wins a test chosen in advance.
\iqlookfor{what each measures and risks.}
\end{solution}

\begin{solution}{iq:ml:monitoring-and-retraining:5}
Not first. Check whether the drift is a data error (units, schema, a vendor change) or a real change in the population; fix data
errors at the source. If the population really changed, check the outcomes: a model can be robust to \omterm{def:ml:online-learning-and-drift:drift}{covariate shift}; retrain
when the outcomes or a revalidation say it is not, and promote through shadow.
\iqlookfor{data path first, outcomes before retraining.}
\end{solution}

\begin{solution}{iq:ml:monitoring-and-retraining:6}
Measure each monitor's false-page rate and its delay on the failures it exists for; drop monitors that never beat a blind one;
recalibrate the rest to a rate the team will answer (count pages at onset); route each alarm to its owner; group related alarms;
and cover slow losses with scheduled reviews instead of pages.
\iqlookfor{calibration to a rate, measured value per monitor, routing.}
\end{solution}
